\section*{第 \ref{ch:b3:residues} 章 --- 洛朗级数与
留数定理}

\begin{solution}{exo:b3:residues:1}
$\frac{\sin z}z = 1 - \frac{z^2}6 + \cdots$：可去奇点，
留数为 $0$。$\frac{\eu^z - 1}{z^2} = \frac1z + \frac12 +
\frac z6 + \cdots$：一级极点，留数为 $1$。
$\frac1{z(z-1)^2}$：$0$ 处的一级极点，留数为
$\frac1{(0-1)^2} = 1$；$1$ 处的二级极点，留数为
$\frac{\dd}{\dd z}\bigl(\frac1z\bigr)\big|_{z=1} = -1$.
$\frac{\cos z}{z^3} = \frac1{z^3} - \frac1{2z} + \cdots$：极点，
阶数为 $3$，留数为 $-\frac12$。$\eu^{1/z} = \sum_{n\geq0}
\frac{z^{-n}}{n!}$：本性奇点，留数为 $1$。$\frac1{\sin z}$：
$n\pi$ 处为一级极点；留数 $\frac1{\cos 0} = 1$，位于 $0$ 处，
留数 $\frac1{\cos\pi} = -1$，位于 $\pi$ 处
（见\cref{met:b3:residues:compute}，$g/h'$）。
\end{solution}

\begin{solution}{exo:b3:residues:2}
令 $z = \eu^{\iu t}$，$\cos t = \frac{z + z^{-1}}2$，$\dd t
= \frac{\dd z}{\iu z}$:
\[
\int_0^{2\pi}\frac{\dd t}{a + \cos t}
= \oint_{\abs z = 1}\frac{2\,\dd z}{\iu\,(z^2 + 2az + 1)} .
\]
根 $z_\pm = -a \pm \sqrt{a^2 - 1}$ 满足 $z_+z_- =
1$，且 $\abs{z_+} < 1 < \abs{z_-}$；$z_+$ 处的留数为
$\frac{1}{z_+ - z_-} = \frac1{2\sqrt{a^2-1}}$，于是积分
值为 $\frac2\iu\cdot2\iu\pi\cdot\frac1{2\sqrt{a^2-1}} =
\frac{2\pi}{\sqrt{a^2-1}}$。当 $a \to 1^+$ 时积分趋于无穷（被积函数在
$t = \pi$ 处达到峰值）；当 $a \to \infty$ 时，积分的行为
类似 $\frac{2\pi}a$，与 $\int\frac{\dd t}a$ 相符。
\end{solution}

\begin{solution}{exo:b3:residues:3}
第一个积分：$\frac{z^2}{1+z^6}$ 在上半平面的极点为 $p =
\eu^{\iu\pi/6}, \iu, \eu^{5\iu\pi/6}$；在每个极点处，
$\operatorname{Res} = \frac{p^2}{6p^5} = \frac{p^3}{6p^6} =
-\frac{p^3}6$，且 $p^3$ 取值 $\iu, -\iu, \iu$：
留数之和为 $-\frac{\iu}{6}$。圆弧部分的贡献为
$O(R^{-4})\cdot O(R) \to 0$:
\[
\int_\R\frac{x^2\,\dd x}{1 + x^6} =
2\iu\pi\Bigl(-\frac \iu6\Bigr) = \frac\pi3 .
\]
第二个积分：$\iu$ 处为二级极点，$\frac1{(1+z^2)^2} =
\frac1{(z-\iu)^2(z+\iu)^2}$:
\[
\operatorname{Res} = \frac{\dd}{\dd z}\,(z +
\iu)^{-2}\Big|_{z=\iu} = \frac{-2}{(2\iu)^3} =
\frac{-2}{-8\iu} = -\frac\iu4,
\qquad
\int_\R\frac{\dd x}{(1+x^2)^2} =
2\iu\pi\cdot\Bigl(-\frac\iu4\Bigr) = \frac\pi2 ,
\]
与 \cref{exo:b3:fouriertransform:3}(b) 一致。
\end{solution}

\begin{solution}{exo:b3:residues:4}
在上半平面内对 $\frac{\eu^{\iu tz}}{z^2 + a^2}$ 取闭合围道
（$t \geq 0$）。此时 $\abs{\eu^{\iu tz}} =
\eu^{-t\operatorname{Im}z} \leq 1$，故圆弧的贡献为
$O(R^{-2})\cdot O(R) \to 0$。被围在内的唯一极点 $\iu a$
为一级极点，留数为 $\frac{\eu^{-at}}{2\iu a}$：
\[
\int_\R\frac{\eu^{\iu tx}}{x^2 + a^2}\dd x =
\frac{\pi}{a}\,\eu^{-at},
\qquad\text{因此}\qquad
\int_\R\frac{\cos(tx)}{x^2+a^2}\dd x = \frac\pi
a\,\eu^{-a\abs t}
\]
（取实部；关于 $t$ 为偶函数）。此式是
$\widehat{\eu^{-a\abs x}} = \frac{2a}{a^2+\xi^2}$ 的反演对应
（见\cref{exo:b3:fouriertransform:1}）：两种计算
相互印证，借助
\cref{thm:b3:fouriertransform:inversion}.
\end{solution}

\begin{solution}{exo:b3:residues:5}
部分分式：$f = \frac1{z-2} - \frac1{z-1}$。
在 $\abs z < 1$ 内（泰勒展开）：
$f = \sum_{n\geq0}\bigl(1 - 2^{-n-1}\bigr)z^n$.
在 $1 < \abs z < 2$ 内：$\frac1{z-2} =
-\sum_{n\geq0}\frac{z^n}{2^{n+1}}$，且 $-\frac1{z-1} =
-\sum_{n\geq1}z^{-n}$：这是一个真正的双边级数。
在 $\abs z > 2$ 内：$f = \sum_{n\geq1}\bigl(2^{n-1} -
1\bigr)z^{-n}$。三者不同，原因在于洛朗展开
依附于\emph{圆环域}，而非点：每个区域都有
自身的几何级数展开。$z^{-1}$ 项的系数（分别为 $-1$ 与
$0$）分别是绕
\emph{内部奇点}的圆周积分，而不是 $0$ 处的留数（$f$ 在
$0$ 处全纯）：对 $1 < \abs z < 2$，系数
为 $-1$，即 $\operatorname{Res}(f, 1)$；对 $\abs z > 2$，系数
为 $0$，即 $\operatorname{Res}(f,1) +
\operatorname{Res}(f,2) = -1 + 1$。$f$ 的留数为 $-1$，
位于 $1$ 处，$+1$ 位于 $2$ 处。
\end{solution}

\begin{solution}{exo:b3:residues:6}
(a) 洛朗级数 $\sum_nz^{-n}/n!$ 含有无穷多个
负幂项，故为本性奇点。解方程 $\eu^{1/z} = w$（$w \neq
0$）：$\frac1z = \log\abs w + \iu\arg w + 2\iu\pi k$，于是
\[
z_k = \frac1{\log\abs w + \iu\arg w + 2\iu\pi k}
\xrightarrow[k\to\infty]{} 0 :
\]
每个非零值在 $0$ 附近都被取到无穷多次 ---
比稠密性更强。
(b) 沿 $z = 1/x$，$x \to +\infty$：$\eu^x \to \infty$；
沿 $z = \iu/y$：模为 $1$。极点要求
$\abs{f(z)} \to \infty$，\emph{沿任意趋近方式}：此处
极限根本不存在，即使在 $[0, +\infty]$ 中也是如此。
\end{solution}

\begin{solution}{exo:b3:residues:7}
(a) 在 $\abs z = 1$ 上：$\abs{z^7 + z - 1} \leq 3 < 4 =
\abs{-4z^3}$。取 $f = -4z^3$，$g = z^7 + z - 1$，用鲁歇定理：
圆盘内有三个零点。
(b) 在 $\abs z = 1$ 上：$\abs{z^4 + 1} \leq 2 < 5 = \abs{5z}$，
$\abs z < 1$ 内有一个零点。在 $\abs z = 2$ 上：$\abs{5z + 1} \leq
11 < 16 = \abs{z^4}$：$\abs z < 2$ 内有四个零点，因此
环域内恰有三个零点。
(c) 设 $P = z^n + a_{n-1}z^{n-1} + \dots$，在 $\abs z = R >
1 + \sum\abs{a_k}$ 上，$\abs{P - z^n} \leq
\bigl(\sum\abs{a_k}\bigr)R^{n-1} < R^n = \abs{z^n}$，于是 $P$
恰有 $n$ 个零点，位于 $D(0, R)$ 内 --- 计重数的达朗贝尔–高斯定理，
仅凭计数即得。
\end{solution}

\begin{solution}{exo:b3:residues:8}
(a) 设 $f(a) = 0$，$f \not\equiv 0$：取 $r$ 使 $f$
在圆周 $C = \partial D(a, r)$ 上无零点（零点孤立），
并令 $m = \min_C\abs f > 0$，由
\cref{thm:b3:holomorphic:weierstrassconv}，$f_n \to f$ 且
$f_n' \to f'$ 在 $C$ 上一致收敛；对充分大的 $n$，$\abs{f_n}
\geq m/2$ 在 $C$ 上，于是
\[
\frac1{2\iu\pi}\int_C\frac{f_n'}{f_n}
\longrightarrow \frac1{2\iu\pi}\int_C\frac{f'}{f} \geq 1
\]
（极限计入零点 $a$；收敛是因为
分子一致收敛，分母一致
有正下界）。左边是一个整数，它计数
$f_n$ 在圆盘内的零点个数，最终须 $\geq 1$ ——
与无零点矛盾。故 $f$ 无零点。

(b) 固定 $w \in \Omega$，在连通开集
$\Omega\setminus\{w\}$（从 $\C$ 的开连通子集中去掉一点，连通性保持）上应用 (a)，
将其用于 $g_n(z) = f_n(z) -
f_n(w)$，它在其上无零点（由单射性），且收敛于 $g = f
- f(w)$。若 $f$ 非常值，则 $g \not\equiv 0$ 在
$\Omega\setminus\{w\}$ 上，故 $g$ 在其上无零点：$f(z) \neq
f(w)$ 对一切 $z \neq w$ 成立。由于 $w$ 任意，$f$ 为
单射。
\end{solution}

\begin{solution}{exo:b3:residues:9}
扇形的边界由 $[0, R]$、弧 $A_R$ 及
射线 $\eu^{2\iu\pi/n}[0, R]$（取反向）组成。内部仅含
一个极点 $p = \eu^{\iu\pi/n}$（即 $\frac1{1+z^n}$ 的极点），其
留数为 $\frac1{np^{n-1}} = \frac{p}{np^n} = -\frac pn$。在
返回射线上，$z = \eu^{2\iu\pi/n}x$ 给出 $z^n = x^n$ 及
$\dd z = \eu^{2\iu\pi/n}\dd x$；圆弧部分为
$O(R^{-n})\cdot O(R) \to 0$。因此
\[
\bigl(1 - \eu^{2\iu\pi/n}\bigr)
\int_0^\infty\frac{\dd x}{1 + x^n}
= 2\iu\pi\Bigl(-\frac{\eu^{\iu\pi/n}}n\Bigr),
\quad\text{故}\quad
\int_0^\infty\frac{\dd x}{1+x^n}
= \frac{2\iu\pi}{n\,\bigl(\eu^{\iu\pi/n} -
\eu^{-\iu\pi/n}\bigr)} = \frac{\pi}{n\sin(\pi/n)} .
\]
$n = 2$ 时：$\frac\pi{2\sin(\pi/2)} = \frac\pi2 =
[\arctan]_0^\infty$。当 $n \to \infty$ 时，该值趋于
$1$，且被积函数趋于 $\mathbf 1_{\intco01}$
（控制收敛定理的控制函数取 $\min(1, x^{-2})$，对 $n \geq 2$）。
\end{solution}

\begin{solution}{exo:b3:residues:10}
在 $\abs z = R$ 充分大时，$\abs f \leq C R^{-2}$：$\abs{\oint_{C_R}
f} \leq 2\pi R\cdot CR^{-2} \to 0$。但对 $R$ 超过全部
极点，留数定理给出 $\oint_{C_R}f =
2\iu\pi\sum_{\text{所有 }p}\operatorname{Res}(f, p)$：全部留数之和
为零。对 $f = \frac1{z(z-1)(z-2)}$，留数为
$\frac1{(-1)(-2)} = \frac12$ 在 $0$ 处，$\frac1{1\cdot(-1)} = -1$
在 $1$ 处，$\frac1{2\cdot1} = \frac12$ 在 $2$ 处 --- 和为 $0$，
与预期一致，并且
\[
\frac1{z(z-1)(z-2)} = \frac{1/2}{z} - \frac1{z - 1} +
\frac{1/2}{z-2} :
\]
留数\emph{恰为}部分分式的系数，
零和恒等式提供了一个免费的一致性检验（或
由其余系数确定最后一个系数）。
\end{solution}

\begin{solution}{exo:b3:residues:11}
在键孔围道上，取定分支如下：
在割线上方 $\log z = \ln x$；在下方，$\log z =
\ln x + 2\iu\pi$。四段之和给出
\[
\Bigl(1 - \eu^{2\iu\pi(a-1)}\Bigr)\int_\varepsilon^R
\frac{x^{a-1}}{1+x}\dd x + \int_{C_R} + \int_{C_\varepsilon}
= 2\iu\pi\operatorname{Res}(f, -1) .
\]
圆弧部分：$\abs{f} \leq \frac{R^{a-1}}{R - 1}$ 在 $C_R$ 上，长为
$2\pi R$，贡献 $O(R^{a-1}) \to 0$（$a < 1$）；
$\abs f \leq \frac{\varepsilon^{a-1}}{1 - \varepsilon}$ 在
$C_\varepsilon$ 上，长 $2\pi\varepsilon$，贡献为 $O(\varepsilon^a)
\to 0$（$a > 0$）。留数在 $z = -1 = \eu^{\iu\pi}$ 处，
$\operatorname{Res} = \eu^{(a-1)\iu\pi}$。于是
\[
I(a) = \frac{2\iu\pi\,\eu^{\iu\pi(a-1)}}{1 -
\eu^{2\iu\pi(a-1)}}
= \frac{2\iu\pi}{\eu^{-\iu\pi(a-1)} - \eu^{\iu\pi(a-1)}}
= \frac{\pi}{-\sin(\pi(a-1))} = \frac{\pi}{\sin\pi a} .
\]
伽马函数的余元公式：$B(a, 1-a) =
\int_0^1t^{a-1}(1-t)^{-a}\dd t$；作代换 $t =
\frac x{1+x}$，$1 - t = \frac1{1+x}$，$\dd t =
\frac{\dd x}{(1+x)^2}$，则它化为 $\int_0^\infty
\frac{x^{a-1}}{1+x}\dd x = I(a)$，由欧拉公式
$B(a, 1-a) = \Gamma(a)\Gamma(1-a)/\Gamma(1) $
（\cref{pb:b3:lebesgue:1}）给出
$\Gamma(a)\Gamma(1-a) = \frac\pi{\sin\pi a}$ --- 特别地，
$\Gamma(\tfrac12) = \sqrt\pi$ 再次成立。

\end{solution}

\begin{solution}{exo:b3:residues:12}
(a) 在 $\abs z = 1$ 上：$\abs{z^4} = 1 < 7 \leq \abs{8z + 1}$
（$\abs{8z} - 1 = 7$）：$P$ 在 $\mathbb D$ 内的零点数与
$8z + 1$ 相同，即\emph{一个}（位于 $-\frac18$）。

(b) 在 $\abs z = 3$ 上：$\abs{8z + 1} \leq 25 < 81 =
\abs{z^4}$：由与 $z^4$ 比较的鲁歇定理，可知全部四个零点位于
$\abs z < 3$ 内，因而 $4 - 1 = 3$ 个零点位于圆环 $1 <
\abs z < 3$ 中。锐化：若某零点满足 $\abs z = r \geq 2.1$，
则应有 $r^4 = \abs{8z + 1} \leq 8r + 1$，但 $r^4 -
8r - 1$ 在 $r \geq 2$ 时递增，且其值为 $19.45 -
16.8 - 1 = 1.65 > 0$（取 $r = 2.1$），故不可能。因此三个外侧零点
位于 $1 < \abs z < 2.1$ 内。（数值上：一个实零点
位于 $-1.95$ 附近，一对共轭复零点在 $1.04 \pm
1.73\iu$ 附近，模为 $2.02$ ——这正是在半径恰为
$2$ 时鲁歇定理的尝试必然失败的原因：该定理要求严格占优，
而零点恰好位于外侧。）

(c) 在 $\iu\R$ 上无零点：$\operatorname{Re}P(\iu t) = t^4 +
1 \geq 1$。右半平面中的零点：用辐角原理，
沿半圆盘 $\{\abs z \leq
R,\ \operatorname{Re}z \geq 0\}$ 的边界。在大圆弧上，$\arg P
\approx \arg z^4$ 的变化量为 $4\cdot\pi = 2\pi\cdot2$（圆弧所张的
角为 $\pi$）。沿虚轴自
$\iu R$ 下行至 $-\iu R$：$P(\iu t) = (t^4 + 1) + 8\iu t$
始终位于右半平面（$\operatorname{Re} > 0$），故
$\arg P$ 在 $\intoo{-\pi/2}{\pi/2}$ 内变化，并
以净变化 $\to 0$（$R \to \infty$ 时）回到原值（两端点均
$\approx \arg t^4 = 0$）。总环绕数为 $\frac{4\pi + 0}
{2\pi} = 2$：右半平面内恰有\emph{两个}零点 ——
与数值结果一致：共轭对 $\approx
1.04 \pm 1.73\iu$ 的实部为正，而实零点
$\approx -0.125$ 与 $\approx -1.96$ 均为负。
\end{solution}

\begin{solution}{pb:b3:residues:1}
\textbf{1.} $\sin(\pi z)$ 在 $\Z$ 上恰有单零点
（$\sin\pi z = 0$ 当且仅当 $z \in \Z$，且 $(\sin\pi z)' = \pi\cos\pi
z \neq 0$ 于这些点成立），且 $\cos(\pi n) \neq 0$：$\pi\cot(\pi z) =
\pi\cos(\pi z)/\sin(\pi z)$ 在 $\Z$ 上有单极点，其留数为
\[
\operatorname{Res}(\pi\cot\pi z,\ n)
= \frac{\pi\cos(\pi n)}{\pi\cos(\pi n)} = 1
\]
（$g/h'$ 法则，见\cref{met:b3:residues:compute}）。

\textbf{2.} $\frac{\sin(\pi z)}{\pi z} = 1 - \frac{(\pi z)^2}6
+ \frac{(\pi z)^4}{120} - \cdots$ 在
$0$ 附近全纯且不为零，其倒数亦全纯
（商的情形，见\cref{def:b3:holomorphic:holo}），其级数为 $1 +
\frac{(\pi z)^2}{6} + \frac{7(\pi z)^4}{360} + \cdots$
（由 $(1 - u)^{-1}$ 型的系数识别，取 $u =
\frac{(\pi z)^2}6 - \frac{(\pi z)^4}{120}$：$z^4$ 项的
系数为 $\frac1{36} - \frac1{120} = \frac{7}{360}$）。
乘以 $\cos(\pi z) = 1 - \frac{(\pi z)^2}2 + \frac{(\pi
z)^4}{24} - \cdots$，再除以 $z$：
\[
\pi\cot(\pi z) = \frac1z\Bigl[1 + \pi^2z^2\Bigl(\frac16 -
\frac12\Bigr) + \pi^4z^4\Bigl(\frac7{360} - \frac1{12} +
\frac1{24}\Bigr)\Bigr] + \cdots
= \frac1z - \frac{\pi^2}3\,z - \frac{\pi^4}{45}\,z^3 -
\cdots
\]
（$\frac7{360} - \frac{30}{360} + \frac{15}{360} =
-\frac8{360} = -\frac1{45}$）。

\textbf{3.} 竖直边 $z = \pm(N + \frac12) + \iu y$：
由以 $\pi$ 为周期的 $\cot$，$\cot(\pi z) = \cot(\pm\frac\pi2
+ \iu\pi y) = -\tan(\iu\pi y) = -\iu\tanh(\pi y)$，模为
$\leq 1$。水平边 $z = x \pm \iu(N + \frac12)$：由
$\abs{\cot(a + \iu b)}^2 = \frac{\cos^2a +
\sinh^2b}{\sin^2a + \sinh^2b} \leq \frac{1 +
\sinh^2b}{\sinh^2b} = \coth^2 b$，于是
\[
\abs{\cot(\pi z)} \leq \coth\bigl(\pi(N + \tfrac12)\bigr)
\leq \coth(\pi/2) < 1.1 .
\]
两个界均 $\leq 2$。

\textbf{4.} 将\cref{thm:b3:residues:residue}应用于 $F(z) =
\pi\cot(\pi z)f(z)$ 在 $C_N$ 上（$N$ 大于 $f$ 的所有极点）：
\[
\frac1{2\iu\pi}\oint_{C_N}F = \sum_{n=-N}^{N}f(n) +
\sum_p\operatorname{Res}(F, p),
\]
整数极点贡献 $f(n)$（第 1 问；$f$
在该处全纯）。在 $C_N$ 上：$\abs F \leq 2\pi\cdot
C\abs z^{-2} \leq 2\pi C N^{-2}$，且周长为 $8(N +
\frac12)$：积分为 $O(1/N) \to 0$。令 $N \to
\infty$，即得所示的求和公式。

\textbf{5.} 衰减性 $\abs f = O(\abs z^{-2})$ 既消去了
围道积分，\emph{又}使 $\sum\abs{f(n)}$ 收敛。
对 $f(z) = \frac1{z + \frac12}$，对称和
$\sum_{-N}^N\frac1{n + \frac12}$ 裂项相消为 $0$（项
$n$ 与 $-n - 1$ 相消），而右端为
$-\operatorname{Res}\bigl(\frac{\pi\cot\pi z}{z + \frac12},
-\frac12\bigr) = -\pi\cot(-\frac\pi2) = 0$：两者一致 ——
但 $\sum\abs{f(n)}$ 发散；该方法计算的是
对称（主值）极限。

\textbf{6.} $g(z) = \frac{\pi\cot(\pi z)}{z^2}$：极点位于
非零整数处，留数为 $\frac1{n^2}$，而在 $0$
处，由第 2 问，
\[
g(z) = \frac1{z^3} - \frac{\pi^2}{3z} - \frac{\pi^4}{45}z -
\cdots :
\qquad \operatorname{Res}(g, 0) = -\frac{\pi^2}3 .
\]
沿 $C_N$ 的围道积分趋于 $0$，与第 4 问相同
（$\abs{g} = O(N^{-2})$ 在 $C_N$ 上）。于是 $0 =
\sum_{n\neq0}\frac1{n^2} - \frac{\pi^2}3$：$2\zeta(2) =
\frac{\pi^2}3$，$\zeta(2) = \frac{\pi^2}6$。

\textbf{7.} $g(z) = \frac{\pi\cot(\pi z)}{z^4} = \frac1{z^5}
- \frac{\pi^2}{3z^3} - \frac{\pi^4}{45z} - \cdots$：留数
在 $0$ 处等于 $-\frac{\pi^4}{45}$，且 $0 =
2\zeta(4) - \frac{\pi^4}{45}$：$\zeta(4) =
\frac{\pi^4}{90}$。

\textbf{8.} 取 $g = \pi\cot(\pi z)/z^{2k}$，
$0$ 处的留数为 $a_{2k-1}$，即 $z^{2k-1}$ 在
$\pi\cot(\pi z)$ 展开式中的系数，围道积分趋于零，给出
$2\zeta(2k) + a_{2k-1} = 0$：$\zeta(2k) = -a_{2k-1}/2$，
是 $\pi^{2k}$ 的有理倍数，因余切函数的
系数皆为有理数。再做一步除法，得 $a_5 =
-\frac{2\pi^6}{945}$，从而 $\zeta(6) = \frac{\pi^6}{945}$。
对于 $\zeta(3)$，取自然核 $g = \pi\cot(\pi z)/z^3$
给出 $\sum_{n\neq0}\frac1{n^3} = 0$（由于奇性）--- 该
方法只得 $0 = 0$，在结构上对奇数阶 zeta 值视而不见
（$\zeta(3)$ 尚无已知闭形式；其
无理性由阿佩里于 1978 年证明，需要
全然不同的思想）。

\textbf{9.} $F(z) = \frac{\pi\cot(\pi z)}{(z - w)(z + w)}$
在整数处有极点（留数为 $\frac1{n^2 - w^2}$，
注意符号：$f(n) = \frac1{(n-w)(n+w)} = \frac1{n^2 -
w^2}$），并在 $\pm w$ 处有单极点，留数为
$\frac{\pi\cot(\pm\pi w)}{\pm2w} = \frac{\pi\cot(\pi
w)}{2w}$ 各（$\cot$ 为奇函数）。第 4 问的论证
（$\abs f = O(\abs z^{-2})$）给出
\[
\sum_{n\in\Z}\frac1{n^2 - w^2} + \frac{\pi\cot(\pi w)}{w} =
0,
\qquad\text{即}\qquad
\pi\cot(\pi w) = \frac1w + \sum_{n\geq1}\frac{2w}{w^2 -
n^2},
\]
（$n = 0$ 项为 $-\frac1{w^2}$，将 $\pm n$ 两项合并）。正规
收敛在 $\C\setminus\Z$ 的紧集上：当 $\abs w \leq
R$ 且 $n \geq 2R$ 时，$\abs{\frac{2w}{w^2 - n^2}} \leq
\frac{2R}{n^2 - R^2} \leq \frac{8R}{3n^2}$。

\textbf{10.} 对于 $\abs w \leq r < 1$：$\frac{2w}{w^2 - n^2} =
-\frac{2w}{n^2}\cdot\frac1{1 - w^2/n^2} =
-2\sum_{k\geq0}\frac{w^{2k+1}}{n^{2k+2}}$，其中
$\abs{\text{各项}} \leq 2r^{2k+1}/n^{2k+2}$，关于
$(n, k)$ 可和：级数的富比尼定理给出如下重排
\[
\pi\cot(\pi w) = \frac1w -
2\sum_{k\geq0}\zeta(2k+2)\,w^{2k+1} .
\]
与第 2 问比较：$-2\zeta(2) = -\frac{\pi^2}3$，而
$-2\zeta(4) = -\frac{\pi^4}{45}$ --- 二者相同。通往
$\frac{\pi^2}6$ 的途径共三条：帕塞瓦尔恒等式（傅里叶级数）、
弦的格林算子之迹，以及余切函数的两种
展开式；频率上的和、算子之迹与
围道积分三者一致，并非偶然 --- 它们各为同一谱恒等式的
一个侧面。

\textbf{11.} 固定 $R \geq 1$，令 $n_0$ 为
$\geq 2R$ 的最小整数。对于 $\abs z \leq R$ 及 $n \geq n_0$：
$\abs{z^2/n^2} \leq \frac14$，故 $1 - z^2/n^2 \in
\bar D(1, \frac14)$，其中主值对数
全纯，且
\[
\Bigl|\log\Bigl(1 - \frac{z^2}{n^2}\Bigr)\Bigr|
\leq \sum_{k\geq1}\frac1k\,\Bigl|\frac{z^2}{n^2}\Bigr|^k
\leq \frac{\abs{z^2/n^2}}{1 - \abs{z^2/n^2}}
\leq \frac{2R^2}{n^2} :
\]
和 $S(z) = \sum_{n\geq n_0}\log(1 - z^2/n^2)$ 在
$\bar D(0, R)$ 上正规收敛，其部分和
$S_N$ 全纯，且一致地 $\abs{S_N} \leq 2R^2\zeta(2)$。由于
$\abs{\eu^a - \eu^b} \leq
\eu^{\max(\abs a,\abs b)}\abs{a - b}$（线段上的中值估计
），尾部乘积 $\prod_{n_0\leq n\leq N} =
\eu^{S_N}$ 在 $\bar D(0, R)$ 上一致收敛到
无零点的 $\eu^S$。乘以固定多项式
$z\prod_{n<n_0}(1 - z^2/n^2)$：$P_N \to P$ 在
$\bar D(0, R)$ 上一致收敛，且
由\cref{thm:b3:holomorphic:weierstrassconv}，$P$
在此全纯；$R$ 任意，故 $P$ 为整函数。在
$\bar D(0, R)$ 上 $P$ 的零点恰为多项式
前因子的零点 --- 即模 $\leq R$ 的整数，各为单零点
（$(1 - z/n)(1 + z/n)$ 有互异单零点，$\eu^S$
无零点）：$P$ 的零点集为 $\Z$，全为单零点。

\textbf{12.} 对有限乘积作对数求导，
在其零点之外：
\[
\frac{P_N'(z)}{P_N(z)} = \frac1z +
\sum_{n=1}^{N}\frac{-2z/n^2}{1 - z^2/n^2}
= \frac1z + \sum_{n=1}^{N}\frac{2z}{z^2 - n^2} .
\]
在紧集 $K \subseteq \C\setminus\Z$ 上：$P_N \to P$ 且
$P_N' \to P'$ 一致收敛
（见\cref{thm:b3:holomorphic:weierstrassconv}），又 $\min_K
\abs P > 0$（$P$ 仅在 $\Z$ 上为零），故
从某项起 $\abs{P_N} \geq \frac12\min_K\abs P$，且 $P_N'/P_N \to
P'/P$ 在 $K$ 上一致收敛。中间一项收敛于
$\frac1z + \sum_{n\geq1}\frac{2z}{z^2-n^2} = \pi\cot(\pi z)$
由第 9 问，故 $P'/P = \pi\cot(\pi z)$ 在
$\C\setminus\Z$ 上成立。

\textbf{13.} $\sin(\pi z)$ 与 $P$ 均为整函数，且有相同
的零点集 $\Z$，零点皆为单零点（第 1 问与第 11 问）。在
$m \in \Z$ 附近，记 $\sin(\pi z) = (z - m)\,\sigma(z)$，且
$P(z) = (z - m)\,\psi(z)$，其中 $\sigma, \psi$ 全纯
且在 $m$ 处不为零（将幂级数分解）：$Q =
\sin(\pi z)/P = \sigma/\psi$ 可全纯延拓，且
在每个整数处亦无零点，且在
$\C\setminus\Z$ 上亦无零点，为无零点函数之商。在该处
\[
\frac{Q'}{Q} = \frac{(\sin\pi z)'}{\sin\pi z} -
\frac{P'}{P} = \pi\cot(\pi z) - \pi\cot(\pi z) = 0 ,
\]
故整函数 $Q'$ 在 $\C\setminus\Z$ 上为零，
由连续性知处处为零，于是 $Q$ 为常数。当 $z \to
0$ 时，$\sin(\pi z)/z \to \pi$ 且 $P(z)/z \to 1$，故 $Q =
\pi$：
\[
\sin(\pi z) = \pi z\prod_{n\geq1}
\Bigl(1 - \frac{z^2}{n^2}\Bigr) .
\]

\textbf{14.} 当 $z = \frac12$ 时：$1 = \sin\frac\pi2 =
\frac\pi2\prod_{n\geq1}\bigl(1 - \frac1{4n^2}\bigr) =
\frac\pi2\prod\frac{4n^2-1}{4n^2}$，则
\[
\frac\pi2 = \prod_{n\geq1}\frac{4n^2}{4n^2 - 1}
= \lim_{N\to\infty}\prod_{n=1}^N
\frac{(2n)(2n)}{(2n-1)(2n+1)}
= \lim_{N\to\infty}
\frac{2\cdot2\cdot4\cdot4\cdots(2N)(2N)}
{1\cdot3\cdot3\cdot5\cdots(2N-1)(2N+1)} :
\]
沃利斯乘积，是欧拉因式分解
的一个直接推论。

\textbf{15.} 对 $\abs z \leq r < 1$，每个因子都落在
$D(1, r^2) \subseteq D(1, 1)$，故 $P(z)/z = \exp\bigl(
\sum_{n\geq1}\log(1 - z^2/n^2)\bigr)$：每个部分积
恰为某个部分和的指数，两边同时
取极限（由 $\exp$ 的连续性）。双重级数
\[
\sum_{n\geq1}\log\Bigl(1 - \frac{z^2}{n^2}\Bigr)
= -\sum_{n\geq1}\sum_{k\geq1}\frac{z^{2k}}{k\,n^{2k}}
= -\sum_{k\geq1}\frac{\zeta(2k)}k\,z^{2k}
= -\zeta(2)\,z^2 + O(z^4)
\]
由级数的富比尼定理重排：$\sum_{n,k}
\frac{r^{2k}}{kn^{2k}} \leq \sum_k\zeta(2k)r^{2k} \leq
\zeta(2)\frac{r^2}{1-r^2} < \infty$。由此
\[
P(z) = z\,\exp\bigl(-\zeta(2)z^2 + O(z^4)\bigr)
= z - \zeta(2)\,z^3 + O(z^5) ,
\]
而练习 13 将其与 $\sin(\pi z) = \pi z -
\frac{\pi^3}6z^3 + O(z^5)$：$\pi\zeta(2) = \frac{\pi^3}6$，
即 $\zeta(2) = \frac{\pi^2}6$。加法形式
（练习 10）与乘积形式
算出的是同一个数。

\textbf{16.} 在紧集 $K \subseteq \C\setminus\Z$ 上，
部分和 $S_N = \frac1z + \sum_{n\leq N}\bigl(
\frac1{z-n} + \frac1{z+n}\bigr)$（练习 9，各项
重组为 $\frac{2z}{z^2-n^2} = \frac1{z-n} +
\frac1{z+n}$）一致收敛到 $\pi\cot(\pi z)$，故
由\cref{thm:b3:holomorphic:weierstrassconv}可知 $S_N' \to
(\pi\cot\pi z)' = -\pi^2/\sin^2(\pi z)$，在 $K$ 上一致收敛。
由于 $S_N' = -\sum_{\abs n\leq N}(z - n)^{-2}$：
\[
\frac{\pi^2}{\sin^2(\pi z)} =
\sum_{n\in\Z}\frac1{(z - n)^2} ,
\]
该级数在 $\C\setminus\Z$ 的紧子集上正规收敛
（通项 $O(n^{-2})$）。

\textbf{17.} 当 $z = \frac12$ 时，左边为 $\pi^2$；
右边，$(\frac12 - n)^2 = \frac{(2n-1)^2}4$，其中 $2n -
1$ 恰好遍历全体奇整数一次，当 $n$
遍历 $\Z$ 时：
\[
\pi^2 = \sum_{n\in\Z}\frac{4}{(2n-1)^2}
= 8\sum_{m\geq0}\frac1{(2m+1)^2},
\qquad
\sum_{m\geq0}\frac1{(2m+1)^2} = \frac{\pi^2}8 .
\]
按奇偶性拆分：$\zeta(2) = \frac{\pi^2}8 +
\sum_{n\geq1}\frac1{(2n)^2} = \frac{\pi^2}8 +
\frac{\zeta(2)}4$，故 $\frac34\zeta(2) = \frac{\pi^2}8$ 且
$\zeta(2) = \frac{\pi^2}6$ 再次成立。

\textbf{18.} 记 $c = \cot\theta$，且 $\cot(2\theta) =
\frac{c^2-1}{2c}$：$\cot\theta - 2\cot(2\theta) = c -
\frac{c^2-1}c = \frac1c = \tan\theta$。由此 $\pi\tan(\pi z)
= \pi\cot(\pi z) - 2\pi\cot(2\pi z)$，并由练习 9 在 $z$
与 $2z$ 处的结论，可得
\[
\pi\cot(\pi z) = \frac1z +
\sum_{n\geq1}\frac{2z}{z^2 - n^2},
\qquad
2\pi\cot(2\pi z) = \frac1z +
\sum_{n\geq1}\frac{8z}{4z^2 - n^2} .
\]
两个级数在每个非极点的固定 $z$ 处绝对收敛，
故其差可任意重新组合：在
第二个级数中，偶数项 $n = 2m$ 给出 $\frac{8z}{4z^2 -
4m^2} = \frac{2z}{z^2 - m^2}$，与第一个级数
全部抵消，剩下
\[
\pi\tan(\pi z) = -\sum_{m\geq0}\frac{8z}{4z^2 - (2m+1)^2}
= \sum_{m\geq0}\frac{8z}{(2m+1)^2 - 4z^2} ,
\]
在避开 $\frac12 + \Z$ 的紧集上正规收敛（通项
$O(m^{-2})$）。

\textbf{19.} 对 $\abs z \leq r < \frac12$：
\[
\frac{8z}{(2m+1)^2 - 4z^2} = \frac{8z}{(2m+1)^2}
\sum_{k\geq0}\Bigl(\frac{4z^2}{(2m+1)^2}\Bigr)^{k},
\]
其中 $\sum_{m,k}8r\,(4r^2)^k(2m+1)^{-2k-2} < \infty$，因为
$4r^2 < 1$。由富比尼定理，将二重和重排为
\[
\pi\tan(\pi z) =
\sum_{k\geq0}8\cdot4^k\,\lambda(2k+2)\,z^{2k+1} .
\]
比较 $\pi\tan(\pi z) = \pi^2z + \frac{\pi^4}3z^3 +
O(z^5)$：$z$ 的系数给出 $8\lambda(2) = \pi^2$
--- 再次用到练习 17 --- 而 $z^3$ 的系数给出
$32\lambda(4) = \frac{\pi^4}3$，即 $\lambda(4) =
\frac{\pi^4}{96}$。去掉偶数分母后，
$\lambda(4) = \zeta(4) - 2^{-4}\zeta(4) =
\frac{15}{16}\zeta(4)$：$\zeta(4) = \frac{16}{15}\cdot
\frac{\pi^4}{96} = \frac{\pi^4}{90}$，与练习 7 一致。

\textbf{20.} $\cot\frac\theta2 - \cot\theta =
\frac{\cos\frac\theta2\,\sin\theta -
\cos\theta\,\sin\frac\theta2}{\sin\frac\theta2\,\sin\theta}
= \frac{\sin\frac\theta2}{\sin\frac\theta2\,\sin\theta} =
\frac1{\sin\theta}$，分子为 $\sin(\theta -
\frac\theta2)$。取 $\theta = \pi z$，则练习 9 在
$\frac z2$ 处的式子为 $\pi\cot\frac{\pi z}2 = \frac2z +
\sum_{n\geq1}\frac{4z}{z^2 - 4n^2}$，于是
\[
\frac{\pi}{\sin(\pi z)}
= \pi\cot\frac{\pi z}2 - \pi\cot(\pi z)
= \frac1z + \sum_{n\geq1}\frac{4z}{z^2 - 4n^2}
- \sum_{n\geq1}\frac{2z}{z^2 - n^2} .
\]
绝对收敛允许按奇偶性重新分组：偶数 $n
= 2m$ 在被减级数中贡献
$\frac{2z}{z^2-4m^2}$，第一个和式则留下 $+\frac{2z}{z^2-4m^2}$，
而奇数 $n$ 以符号 $-$ 保留：
\[
\frac{\pi}{\sin(\pi z)} = \frac1z +
\sum_{n\geq1}(-1)^n\,\frac{2z}{z^2 - n^2} .
\]
符号检验：$\operatorname{Res}\bigl(\pi/\sin(\pi z),
n\bigr) = \pi/(\pi\cos\pi n) = (-1)^n$
（见\cref{met:b3:residues:compute}），且 $(-1)^n\frac{2z}{z^2 -
n^2} = \frac{(-1)^n}{z-n} + \frac{(-1)^n}{z+n}$ 在
$\pm n$ 处恰好具有该留数。

\textbf{21.} 设 $\abs z \leq r < 1$，按
练习 19 的方式展开每一项（$(-1)^n\frac{2z}{z^2-n^2} =
(-1)^{n-1}\frac{2z}{n^2}\sum_k\frac{z^{2k}}{n^{2k}}$），并
应用富比尼定理：
\[
\frac{\pi}{\sin(\pi z)} = \frac1z +
2\sum_{k\geq0}\eta(2k+2)\,z^{2k+1},
\qquad
\eta(s) = \sum_{n\geq1}\frac{(-1)^{n-1}}{n^s} .
\]
泰勒展开一侧：$\sin(\pi z) = \pi z(1 - \frac{(\pi z)^2}6 +
O(z^4))$ 给出 $\frac\pi{\sin\pi z} = \frac1z +
\frac{\pi^2}6z + O(z^3)$：$2\eta(2) = \frac{\pi^2}6$，于是
$\eta(2) = \frac{\pi^2}{12}$。一致性检验：$\eta(2) =
\zeta(2) - 2\sum_n(2n)^{-2} = (1 - 2^{1-2})\zeta(2) =
\frac{\zeta(2)}2 = \frac{\pi^2}{12}$.

\textbf{22.} 取 $z = \frac12$，由练习 20 可得
\[
\pi = 2 + \sum_{n\geq1}(-1)^n\frac{1}{\frac14 - n^2}
= 2 + 4\sum_{n\geq1}\frac{(-1)^{n-1}}{(2n-1)(2n+1)} .
\]
利用 $\frac1{(2n-1)(2n+1)} = \frac12\bigl(\frac1{2n-1} -
\frac1{2n+1}\bigr)$，两个交错级数均
收敛（由莱布尼茨判别法），和式分解为 $\frac12\bigl[L -
(1 - L)\bigr] = L - \frac12$，其中 $L = 1 - \frac13 +
\frac15 - \cdots$ 且 $\frac13 - \frac15 + \frac17 - \cdots
= 1 - L$。于是 $\pi = 2 + 4(L - \frac12) = 4L$：
\[
\frac\pi4 = 1 - \frac13 + \frac15 - \frac17 + \cdots
\]
小结：每个核都是亚纯函数，其极点位于一个
等差数列上，并带有给定的留数。
余切在每个整数处留数为 $1$，并对
$f(n)$ 求和；其导数使极点平方，并
给出 $\lambda(2)$；正切把极点移至
$\frac12 + \Z$ 处，并给出奇数分母对应的值；$\pi/\sin$ 保留
整数极点，但留数交替为 $(-1)^n$，故
出现交错级数。一阶核全是
奇函数：将 $n$ 与 $-n$ 配对，使偶次幂系数加倍，
奇次幂系数则消为零，故 $\zeta(2k)$
可机械地得出，而 $\zeta(3)$ 从不出现。
这一盲区源于奇偶性，而非技巧上的欠缺。

\textbf{23.} 取 $w = 2\iu\pi z$：
\[
\iu\pi z + \frac{2\iu\pi z}{\eu^{2\iu\pi z} - 1}
= \iu\pi z\,\frac{\eu^{2\iu\pi z} + 1}{\eu^{2\iu\pi z} - 1}
= \iu\pi z\,\frac{\eu^{\iu\pi z} +
\eu^{-\iu\pi z}}{\eu^{\iu\pi z} - \eu^{-\iu\pi z}}
= \pi z\,\frac{\cos\pi z}{\sin\pi z} = \pi z\cot(\pi z) .
\]
由此，对 $w = 2\iu\pi z$ 利用生成函数（并
$B_1 = -\frac12$ 消去 $\iu\pi z$ 项，奇数 $B$
在此之后均为零）：
\[
\pi z\cot(\pi z) = \sum_{k\geq0}\frac{B_{2k}}{(2k)!}
(2\iu\pi z)^{2k}
= 1 + \sum_{k\geq1}(-1)^k\frac{(2\pi)^{2k}B_{2k}}{(2k)!}
z^{2k} .
\]
另一方面，由部分分式展开（第 IV 部分）
得 $\pi z\cot(\pi z) = 1 - 2\sum_{k\geq1}\zeta(2k)z^{2k}$
（逐项展开 $\frac{2z^2}{z^2 - n^2} =
-2\sum_k\frac{z^{2k}}{n^{2k}}$，并对 $n$ 求和，正规
收敛使得在 $\abs z < 1$ 时可交换求和次序）。
比较系数：$-2\zeta(2k) =
(-1)^k\frac{(2\pi)^{2k}B_{2k}}{(2k)!}$，即所给公式。

\textbf{24.} $k = 1$: $\frac{(2\pi)^2}{2\cdot2}\cdot\frac16
= \frac{\pi^2}6$. $k = 2$: $-\frac{(2\pi)^4}{2\cdot24}
\cdot\bigl(-\frac1{30}\bigr) = \frac{16\pi^4}{48\cdot30} =
\frac{\pi^4}{90}$. $k = 3$: $\frac{(2\pi)^6}{2\cdot720}
\cdot\frac1{42} = \frac{64\pi^6}{1440\cdot42} =
\frac{\pi^6}{945}$.

\textbf{25.} 记 $C(z) = \pi z\cot(\pi z) = 1 -
2\sum_{k\geq1}\zeta(2k)z^{2k}$（见练习 23）。直接
对 $C = \pi z\cot(\pi z)$，利用 $(\cot u)'
= -1 - \cot^2u$:
\[
zC'(z) = \pi z\cot(\pi z) - \pi^2z^2\bigl(1 +
\cot^2(\pi z)\bigr) = C - \pi^2z^2 - C^2 .
\]
现将两边按 $z^2$ 的幂展开。左边：
$\sum_k(-4k)\,\zeta(2k)\,z^{2k}$。右边：将
级数平方，
\[
C - C^2 = 2\sum_{k\geq1}\zeta(2k)z^{2k}
- 4\sum_{k\geq2}\Bigl(\sum_{j=1}^{k-1}\zeta(2j)
\zeta(2k-2j)\Bigr)z^{2k},
\]
而项 $-\pi^2z^2 = -6\zeta(2)z^2$ 只影响 $k =
1$ 的情形。比较 $z^{2k}$ 的系数（$k \geq 2$）：
\[
-4k\,\zeta(2k) = 2\,\zeta(2k) -
4\sum_{j=1}^{k-1}\zeta(2j)\,\zeta(2k-2j),
\]
即 $\bigl(k + \frac12\bigr)\zeta(2k) =
\sum_{j=1}^{k-1}\zeta(2j)\zeta(2k-2j)$。（当 $k = 1$ 时
恒等式为 $-4\zeta(2) = 2\zeta(2) - 6\zeta(2)$：一
个一致性检验，而非新信息）。应用：$k =
2$：$\frac52\zeta(4) = \zeta(2)^2 = \frac{\pi^4}{36}$，故
$\zeta(4) = \frac{\pi^4}{90}$；$k = 3$：$\frac72\zeta(6) =
2\zeta(2)\zeta(4) = \frac{\pi^6}{270}$，故 $\zeta(6) =
\frac{2}{7}\cdot\frac{\pi^6}{270} = \frac{\pi^6}{945}$。只
需一个超越起点（$\zeta(2) = \frac{\pi^2}6$），纯代数便
可生成全部偶数指标的 ζ 值。
\end{solution}
